\documentclass[10pt,twocolumn]{article}
\usepackage[margin=0.55in,top=0.5in,bottom=0.45in]{geometry}
\usepackage{graphicx}
\usepackage{booktabs}
\usepackage{enumitem}
\usepackage{caption}
\usepackage{float}
\usepackage{xcolor}
\usepackage{titlesec}
\usepackage[hidelinks]{hyperref}

\setlength{\parskip}{3pt}
\setlength{\parindent}{0pt}
\setlist{nosep,leftmargin=*}

\titlespacing*{\section}{0pt}{6pt}{2pt}
\titlespacing*{\subsection}{0pt}{4pt}{1pt}
\titleformat{\section}{\normalfont\bfseries\normalsize}{\thesection}{0.5em}{}
\titleformat{\subsection}{\normalfont\bfseries\small}{\thesubsection}{0.5em}{}

\setlength{\abovecaptionskip}{4pt}
\setlength{\belowcaptionskip}{2pt}
\captionsetup{font=small}

\begin{document}

\twocolumn[{
\begin{center}
{\Large\bfseries Autoreason: Autoresearch for Subjective Domains}\\[4pt]
{\normalsize Nous Research}\\[2pt]
{\small\url{https://github.com/NousResearch/autoreason}}\\[6pt]
\end{center}
}]

\section*{Abstract}

We introduce autoreason, a method for iterative LLM refinement in subjective domains where no objective metric exists. Building on autoresearch paradigms, autoreason addresses systematic failures in LLM self-improvement: sycophancy, overcriticism, overcompromise, authorship bias, scope drift, and context collapse. Our method employs fresh isolated agents per role, blind judge panels with ranked choice voting and Borda count aggregation, and conservative tiebreaking. In a single experiment on go-to-market strategy generation, autoreason achieved the highest Borda score (35/35) and all first-place rankings (7/7 judges) against four baseline methods after 15 iterations. The approach ran for 26 passes without reaching the predefined convergence threshold of 3 consecutive incumbent wins, achieving 2 consecutive wins at passes 14--15 and 24--25. We present this as an exploratory study of one approach to subjective iteration, not a validated general method.

\section{Introduction}

Large language models fail systematically when iterating on subjective work. Given a task to improve their own output, LLMs exhibit predictable pathologies: they agree too readily with criticism (sycophancy), find fault where none exists (overcriticism), dilute strong positions into bland compromises (overcompromise), recognize and favor their own prior outputs (authorship bias), expand scope beyond the original task (scope drift), and lose track of quality baselines as context accumulates (context collapse).

These failures persist across prompting strategies. Recent work demonstrates that even in the relatively objective domain of code generation, iterative improvement prompts lead to progressive quality degradation~\cite{slopcodebench}. Common strategies like ``critique and revise,'' ``harsh critic,'' and ``improve this'' all suffer from the same fundamental issue: the model cannot maintain stable quality judgments across iterations.

The problem is particularly acute in subjective domains---strategy documents, creative writing, analysis, and planning---where no ground truth metric exists. Without an objective fitness function, traditional autoresearch approaches cannot be applied directly. Human evaluation remains the gold standard but is expensive and slow.

We present autoreason, a method that constructs a subjective fitness function through independent blind evaluation. By isolating each role into fresh agents and aggregating preferences across judge panels, autoreason achieved iterative improvement in our test case. We emphasize this represents initial exploration on a single task, not a comprehensively validated approach.

\section{Method}

Autoreason implements an A/B/AB evaluation loop where:
\begin{itemize}
\item \textbf{A} is the current incumbent (best version so far)
\item \textbf{B} is a new challenger created to improve upon A
\item \textbf{AB} is a synthesis that attempts to combine the best of both A and B
\end{itemize}

The core insight is that each role must be performed by a fresh, isolated agent with no knowledge of prior iterations or the other agents' identities.

\subsection{Architecture}

Each pass consists of four phases:

\textbf{Strawman Phase}: A fresh critic agent receives the original task and the current incumbent A. The critic identifies potential weaknesses, areas for improvement, or missing elements without seeing any prior criticism or iteration history.

\textbf{Author Phase}: Two fresh author agents work independently. The first receives the original task, incumbent A, and the strawman critique to generate challenger B. The second receives the same inputs plus both A and B to create synthesis AB. Neither author has knowledge of previous iterations or judge feedback. Authors use temperature 0.8 to encourage creative exploration.

\textbf{Judge Phase}: A panel of three fresh judge agents independently evaluates A, B, and AB in randomized order. Judges receive only the original task and the three candidates, with no indication of which is incumbent, challenger, or synthesis. Each judge provides a complete ranking from best to worst. Judges use temperature 0.3 for more consistent evaluation. We use standard Borda count scoring where first place receives 2 points, second place receives 1 point, and third place receives 0 points.

\textbf{Aggregation}: Rankings are aggregated by summing Borda scores across all judges. The candidate with the highest total score becomes the new incumbent. In case of equal total scores across judges, the current incumbent A wins (conservative tiebreak).

\subsection{Convergence}

The system tracks consecutive wins by the incumbent A, resetting the count when A loses. We defined convergence as 3 consecutive A wins, though our experiment never reached this threshold. The system achieved 2 consecutive wins at two points (passes 14--15 and 24--25). The fact that AB won approximately 50\% of passes suggests the system continued finding viable improvements throughout the experiment.

\subsection{Design Rationale}

Each design choice addresses a specific failure mode observed in preliminary experiments:

\textbf{Fresh isolated agents} prevent context accumulation and authorship bias. An author who remembers creating A will be biased when creating B. A judge who knows the iteration history loses calibration.

\textbf{Blind evaluation} prevents judges from favoring the challenger due to implicit ``newer is better'' assumptions or recognizing authorship patterns.

\textbf{Judge panels} reduce single-judge noise and provide more robust preferences through aggregation. We chose 3 judges for the iterative process as a balance between robustness and computational cost, though we acknowledge this choice was not empirically validated.

\textbf{Three candidates (A/B/AB)} allow the system to explore both incremental improvements (B) and synthetic combinations (AB), increasing the chance of finding better solutions.

\textbf{Conservative tiebreak} prevents oscillation between equally good alternatives and provides a stability mechanism.

\section{Experiments}

We evaluated autoreason on go-to-market strategy generation for an open-source Kubernetes CLI tool. This task represents a typical subjective business document where quality is clear to human judges but no objective metric exists.

\subsection{Experimental Setup}

All experiments used claude-sonnet-4-20250514 with temperature 0.8 for authors and 0.3 for judges. The v2 architecture employed 3-judge panels with Borda count aggregation and conservative tiebreaking. Each pass required approximately 6 LLM calls (1 strawman, 2 authors, 3 judges), totaling approximately 156 calls for 26 passes.

\subsection{Main Experiment}

We ran the full autoreason process for 26 passes. The system never reached the convergence threshold of 3 consecutive incumbent wins, achieving at most 2 consecutive incumbent wins at passes 14--15 and 24--25. The win/loss trajectory showed early rapid improvement, word count oscillation in middle passes, and eventual quality plateau.

\subsection{Baseline Comparison}

We compared autoreason's pass 15 output against four baseline methods, each given the same initial task:

\begin{itemize}
\item \textbf{Conservative}: Single-shot generation (no iteration)
\item \textbf{Improve\_this}: Direct iterative improvement prompt
\item \textbf{Harsh\_critic}: Iterative improvement with aggressive criticism
\item \textbf{Critique\_and\_revise}: Structured critique followed by revision
\end{itemize}

A panel of 7 fresh judges evaluated all five outputs using ranked choice voting with standard Borda count scoring (4 points for 1st place through 0 for 5th, maximum 35 points per method). We selected pass 15 as it represented the first occurrence of 2 consecutive incumbent wins.

\subsection{Temporal Comparison}

To test whether continued iteration affected quality, we compared outputs from pass 15 versus pass 25 using a separate 7-judge blind panel.

\subsection{Baseline Visibility Test}

We tested whether judges needed a baseline for drift detection by comparing autoreason output against the critique-and-revise output with and without the original task output shown for reference.

\section{Results}

\subsection{Baseline Comparison}

\begin{table}[H]
\centering
\small
\begin{tabular}{@{}lccr@{}}
\toprule
\textbf{Method} & \textbf{Borda} & \textbf{1st} & \textbf{Words} \\
\midrule
\textbf{Autoreason} & \textbf{35/35} & \textbf{7/7} & 1,800 \\
Conservative & 21 & 0 & 862 \\
Improve\_this & 18 & 0 & 2,116 \\
Harsh\_critic & 18 & 0 & 1,961 \\
Critique \& revise & 13 & 0 & 2,507 \\
\bottomrule
\end{tabular}
\caption{7-judge blind panel results. Maximum Borda score is 35.}
\label{tab:baselines}
\end{table}

Autoreason's pass 15 output achieved the highest Borda score (35/35) and received all 7 first-place votes. The conservative baseline scored 21/35, outperforming all iterative baselines. Critique and revise, the most commonly used iterative approach, scored lowest at 13/35.

Word counts revealed length inflation across methods: from 847 words at pass 1, autoreason reached 1,800 words by pass 15, while improve\_this reached 2,116, harsh\_critic 1,961, and critique\_and\_revise ballooned to 2,507. The conservative baseline remained at 862 words.

\subsection{Quality Trajectory}

Pass 15 beat pass 25 by a 6--1 margin among judges, suggesting that continued iteration beyond the first stable point may not improve quality.

\subsection{Baseline Visibility}

With the original task baseline shown, judges preferred autoreason output 7--0 against the critique-and-revise version. Without baseline shown, the preference was 3--2. Judges need reference points to detect drift effectively.

\subsection{Qualitative Improvements}

The initial output contained generic targets and unsupported metrics (``\$49/user pricing,'' ``\$100K MRR''). Pass 15 output included quantified pain points (``\$15K/incident $\times$ 6/yr''), team pricing (``\$1,499/mo''), customer validation references (``50+ interviews, 75\% pilot success''), and unit economics (``CAC \$2K, LTV \$54K'').

\section{Related Work}

The autoresearch paradigm~\cite{autoresearch} demonstrates success in objective domains with clear metrics. SlopCodeBench~\cite{slopcodebench} shows that iterative refinement fails even for code generation, with prompting strategies unable to prevent quality collapse. Work on context collapse~\cite{ace} identifies how LLMs lose quality baselines with accumulated context. LLM Council~\cite{llmcouncil} explores multi-agent deliberation architectures. Our use of blind evaluation panels builds on established peer review practices rather than representing a novel contribution.

\section{Discussion}

\subsection{Bloat/Prune Oscillation}

Word count variation between passes (ranging from approximately 1,600 to 2,100 words in middle passes) suggests authors alternate between adding detail and simplifying when improvement directions are unclear. This oscillation indicates that the task specification may need refinement rather than more iteration.

\subsection{Convergence Patterns}

Our experiment never reached the predefined convergence threshold of 3 consecutive incumbent wins, achieving at most 2 consecutive wins at passes 14--15 and 24--25. The fact that AB won approximately 50\% of passes suggests the synthesis mechanism successfully found improvements even late in the process. Whether this represents oscillation between equally viable solutions or indicates the method doesn't converge requires testing across multiple tasks.

\subsection{Judge Calibration}

Judges maintain better calibration with baselines. Our baseline visibility test showed unanimous preference (7--0) with baseline against the adversarial version, versus 3--2 without, demonstrating that judges need reference points to detect drift effectively.

\subsection{Applications Beyond Subjective Domains}

Autoreason may prove valuable even in domains with objective metrics. Metrics often fail to capture all relevant quality dimensions---a function might pass all tests yet be unmaintainable. Autoreason could provide a complementary signal in such cases.

\subsection{Limitations}

This work represents an exploratory study of a single task domain (go-to-market strategy) with one experimental run. We did not perform statistical significance testing, analyze inter-rater agreement, or test alternative numbers of candidates beyond our A/B/AB design. All results come from a single model (claude-sonnet-4-20250514) and may not generalize to other LLMs. The method's effectiveness on other subjective tasks remains untested.

\section{Conclusion}

Autoreason demonstrates one approach to iterative refinement in subjective domains. By isolating agents, implementing blind evaluation, and using conservative aggregation, the method achieved improvements in our test case while limiting the length inflation seen in baseline approaches. The key insight---that isolation prevents context contamination---emerged from preliminary experiments with approximately 1,800 LLM calls that revealed issues with positional bias, single judge noise, and shared context contamination. Future work should test across diverse domains, include statistical validation, and analyze optimal judge panel sizes.

\begin{thebibliography}{9}
\small
\bibitem{autoresearch} A. Karpathy. Autoresearch, 2024. \url{https://github.com/karpathy/autoresearch}

\bibitem{slopcodebench} G. Orlanski et al. SlopCodeBench: Benchmarking how coding agents degrade over iterative tasks. \textit{arXiv:2603.24755}, 2026.

\bibitem{ace} J. Zhang et al. ACE: Agentic context engineering for iterative rewriting. \textit{ICLR}, 2026.

\bibitem{llmcouncil} A. Karpathy. LLM Council, 2025. \url{https://github.com/karpathy/llm-council}
\end{thebibliography}

\end{document}
